\section{Preliminaries}
\label{sec:preliminaries}

\paragraph{Terminology} Throughout the paper, we use \textsc{Kyber} and \textsc{ML-KEM} interchangeably, except when distinguishing the original 
\textsc{Kyber} submission from the standardized ML-KEM specification.

\paragraph{Ring and module setup.} Let $n=256$ and $q=3329$, and let $R=\mathbb{Z}[X]/(X^n+1)$ and $R_q=\mathbb{Z}_q[X]/(X^n+1)$ be \textsc{Kyber}'s
base ring, with coefficients centered in $(-q/2,q/2]$. The modulus satisfies $n \mid (q-1)$, needed for the
number-theoretic transform (NTT) that \textsc{Kyber} uses for polynomial multiplication in $R_q$. A module rank $k$ determines the dimensions of the public matrix
$\mathbf{A}\in R_q^{k\times k}$ and the secret/error vectors $\mathbf{s},\mathbf{e}\in R_q^k$. $k=2,3,4$ give ML-KEM-512, -768, and -1024,
respectively~\cite{fips203}.

\paragraph{Centered binomial noise.} \textsc{Kyber} samples secret and error coefficients from the centered binomial distribution $\mathrm{CBD}*\eta$. It draws
$2\eta$ uniform bits $a_1,\dots,a*\eta,b_1,\dots,b_\eta$ and gives $\sum_{i=1}^\eta a_i - \sum_{i=1}^\eta b_i$, a value in ${-\eta,\dots,\eta}$
with variance $\eta/2$. ML-KEM-768 and -1024 use $\mathrm{CBD}*2$ for both the secret and error terms; ML-KEM-512 uses $\mathrm{CBD}*3$ for the secret. We
write $\mathrm{CBD}*\eta(x)$ for its probability mass function. $\mathrm{CBD}*\eta$ is
especially desirable because it can be sampled from raw PRF output in constant time
without rejection and uses only bit-counting and a subtraction.

\paragraph{Module-LWE.} The scheme's IND-CPA security reduces to the decisional Module Learning With Errors (MLWE)
problem (see \cite{DCC:LanSte15}) over $R_q$. An adversary must distinguish
$(\mathbf{A}, \mathbf{A}\mathbf{s}+\mathbf{e})$, where $\mathbf{A}\leftarrow R_q^{k\times k}$ and $\mathbf{s},\mathbf{e}\leftarrow \mathrm{CBD}_\eta^k$, from a uniformly random
pair in $R_q^{k\times k}\times R_q^k$. Increasing $k$ increases the dimension
of the lattice problem and is the main lever \textsc{Kyber} uses to raise its claimed security level~\cite{kyberround1}.


\paragraph{Core-SVP hardness estimates.} We estimate MLWE hardness following the standard core-SVP methodology used throughout the \textsc{Kyber}
submission documents and re-implemented in the \textsc{Kyber} team's public estimator scripts. An adversary is modeled as
running BKZ lattice reduction with block size $b$ against the primal (or dual) lattice attack. The cost of BKZ-$b$ is dominated by a single call to an
SVP oracle in dimension $b$, whose cost is estimated as $2^{cb}$ operations for a constant $c$ (e.g., $c\approx 0.292$ classically and 
$c\approx 0.265$ under a quantum sieve)~\cite{bkz2,estimator}. For a given parameter set, the scripts search over block size 
$b$ and sample count $m$
for the cheapest attack that succeeds against the secret/error distribution's standard deviation, and report the $\log_2$ cost as the 
\emph{classical} and \emph{quantum} core-SVP security level. We report both
throughout. We note that these estimates are conservative lower bounds on attack cost
rather than concrete running times.

\paragraph{Decryption failure.} Because \textsc{Kyber}'s public-key encryption scheme is only approximately correct, decapsulation can fail. The
accumulated LWE and rounding error can push a coefficient of the decrypted message past the threshold at $q/4$ that 
separates a $0$ bit from a $1$ bit. We refer to the decryption-failure probability as the probability, over the randomness of key generation and 
encryption, that at least one of the $n$ message-bit coefficients suffers such an error.
\footnote{Even though a single decryption error may not sound catastrophic, it is well established that once an attacker finds a single failure, 
they can use it as a compass to drastically accelerate finding more failures, a type of attack known as ``directional failure-boosting'' \cite{EC:DAnRosVir20}.}
It is computed by convolving the noise, secret, and 
(when present) compression rounding-error distributions coordinate-wise and taking the tail probability of the final 
distribution beyond $q/4$, scaled by $n$. A failure probability of $2^{-128}$ is considered a standard benchmark which aligns with the target
of 128-bit security. We use the same bound throughout as our correctness target.

\paragraph{From PKE to KEM.} \textsc{Kyber}'s IND-CPA-secure public-key encryption scheme is lifted to an IND-CCA2-secure KEM using a tweaked variant of the 
Fujisaki--Okamoto (FO) transform~\cite{fips203}. The ciphertext produced by encryption consists of two components, $\mathbf{u}\in R_q^k$ and $v\in R_q$, which
ML-KEM serializes after \emph{compression}, where each coefficient is rounded.
This trades a small increase in decryption-failure probability for a smaller ciphertext. The public key,
consisting of $\mathbf{A}$'s seed and $\mathbf{A}\mathbf{s}+\mathbf{e}$, is
never compressed. Throughout the paper, ``ciphertext compression'' refers to
this rounding of $\mathbf{u}$ and $v$. We remove this compression in
\textsc{Kaiburr} to achieve formally verified correctness bounds (see Section~\ref{sec:verification}). The FO
transform itself, the ring, the NTT, and the symmetric primitives are left
untouched.